% ===== PRÉAMBULE CANONIQUE — à coller VERBATIM en tête de CHAQUE .tex =====
\documentclass[11pt,a4paper]{article}
\usepackage[utf8]{inputenc}\usepackage[T1]{fontenc}\usepackage{lmodern}
\usepackage[french]{babel}
\usepackage{amsmath,amssymb,mathtools}\usepackage{icomma}
\usepackage{siunitx}\sisetup{locale=FR}  % \qty{}{}, \si{}, \ang{}
\usepackage{xcolor}\usepackage{colortbl}
\usepackage{tikz}\usepackage{pgfplots}\pgfplotsset{compat=1.18}
\usepackage[siunitx,europeanresistors]{circuitikz} % circuits (élec.)
\usepackage[most]{tcolorbox}\usepackage{enumitem}\usepackage{array,booktabs}
\usepackage[a4paper,margin=2.2cm]{geometry}
\usetikzlibrary{arrows.meta,calc,angles,quotes,positioning,patterns,%
  backgrounds,shapes.geometric,decorations.pathmorphing,decorations.markings,intersections}
\definecolor{primary}{RGB}{222,49,99}\definecolor{primarydark}{RGB}{150,22,55}
\definecolor{defcol}{RGB}{0,114,178}\definecolor{methcol}{RGB}{52,92,148}
\definecolor{excol}{RGB}{0,135,90}\definecolor{attncol}{RGB}{200,40,55}
\tcbset{boxbase/.style={enhanced,breakable,boxrule=0.9pt,arc=2.5pt,
  left=3mm,right=3mm,top=2.3mm,bottom=2.3mm,fonttitle=\bfseries,coltitle=white}}
% --- boîtes TUTORIEL (noms EXACTS, ne pas renommer) ---
\newtcolorbox{cle}[1][]{boxbase,colback=defcol!6,colframe=defcol,
  title={\textsc{À retenir}\ifx\relax#1\relax\relax\else\textmd{ : \itshape #1}\fi}}
\newtcolorbox{methode}[1][]{boxbase,colback=methcol!7,colframe=methcol,sharp corners,
  title={\textsc{Méthode}\ifx\relax#1\relax\relax\else\textmd{ --- \itshape #1}\fi}}
\newtcolorbox{exemplebox}[1][]{boxbase,colback=excol!5,colframe=excol,
  title={\textsc{Exemple}\ifx\relax#1\relax\relax\else\textmd{ : \itshape #1}\fi}}
\newtcolorbox{piege}[1][]{boxbase,colback=attncol!6,colframe=attncol,
  title={\textsc{Piège}\ifx\relax#1\relax\relax\else\textmd{ : \itshape #1}\fi}}
\newtcolorbox{remarque}[1][]{boxbase,colback=black!4,colframe=black!45,
  title={\textsc{Remarque}\ifx\relax#1\relax\relax\else\textmd{ : \itshape #1}\fi}}
% --- boîtes EXERCICE / CORRIGÉ (noms EXACTS, auto-comptées) ---
\newtcolorbox[auto counter]{exo}[1][]{boxbase,colback=primary!4,colframe=primary,
  title={\textsc{Exercice}~\thetcbcounter\ifx\relax#1\relax\relax\else\textmd{ --- \itshape #1}\fi}}
\newtcolorbox[auto counter]{corrige}[1][]{boxbase,colback=excol!4,colframe=excol,
  title={\textsc{Corrigé}~\thetcbcounter\ifx\relax#1\relax\relax\else\textmd{ --- \itshape #1}\fi}}
\newcommand{\vv}[1]{\vec{#1}}\newcommand{\ds}{\displaystyle}
\newcommand{\R}{\mathbb{R}}\newcommand{\N}{\mathbb{N}}\newcommand{\Z}{\mathbb{Z}}
% (un \usepackage{fancyhdr}+en-tête est possible ; il est ignoré côté web)
% =========================================================================
\begin{document}
\sisetup{per-mode=symbol}

\begin{exo}[variations composées de la force]
Deux charges se repoussent avec une force d'intensité $F_0 = \SI{1}{\newton}$.
\begin{enumerate}[label=\alph*)]
  \item On \emph{double} chacune des deux charges \emph{et} on \emph{divise la distance par 2}.
        Que devient l'intensité de la force ?
  \item En repartant de $F_0$, on \emph{divise chaque charge par 3} \emph{et} la \emph{distance par
        3}. Que devient l'intensité ?
\end{enumerate}
\end{exo}

\begin{exo}[deux forces perpendiculaires]
Une charge subit deux forces électriques \textbf{perpendiculaires} entre elles, d'intensités
$F_1 = \SI{6}{\newton}$ et $F_2 = \SI{8}{\newton}$.
\begin{enumerate}[label=\alph*)]
  \item Calcule l'intensité de la force résultante.
  \item Donne l'angle $\theta$ que fait la résultante avec $\vv{F}_1$ ($\tan\theta = F_2/F_1$).
\end{enumerate}
\end{exo}

\begin{exo}[triangle équilatéral]
On place une charge $+q$ avec $q = \SI{2}{\micro\coulomb}$ à chaque sommet d'un triangle
équilatéral de côté $a = \SI{10}{\centi\metre}$.
\begin{center}
\begin{tikzpicture}[scale=0.9,>=Stealth,font=\footnotesize,
    ch/.style={circle,fill=attncol,minimum size=7mm,inner sep=0pt,text=white,font=\bfseries\scriptsize}]
  \coordinate (A) at (-1.5,0); \coordinate (B) at (1.5,0); \coordinate (C) at (0,2.6);
  \draw[black!45] (A)--(B)--(C)--cycle;
  \node[ch] at (A) {$+q$}; \node[ch] at (B) {$+q$}; \node[ch] at (C) {$+q$};
  \node[below,font=\scriptsize] at (0,-0.15) {$a$};
  \draw[->,line width=1pt,defcol] (C) -- ++(122:1.1);
  \draw[->,line width=1pt,defcol] (C) -- ++(58:1.1);
  \draw[->,line width=1.4pt,primary] (C) -- ++(90:1.5) node[above,font=\scriptsize]{$\vv{F}_{\text{tot}}$};
\end{tikzpicture}
\end{center}
\begin{enumerate}[label=\alph*)]
  \item Exprime, par symétrie, l'intensité de la force totale subie par la charge du sommet
        supérieur ($F_{\text{tot}} = 2f\cos 30^\circ = \sqrt{3}\,K q^2/a^2$).
  \item Fais l'application numérique.
\end{enumerate}
\end{exo}

\begin{exo}[trois charges alignées]
Trois charges sont alignées et espacées de $\SI{10}{\centi\metre}$ : $A = \SI{+2}{\micro\coulomb}$,
$B = \SI{+3}{\micro\coulomb}$ (au milieu), $C = \SI{-2}{\micro\coulomb}$.
\begin{center}
\begin{tikzpicture}[scale=1,>=Stealth,font=\footnotesize,
    ch/.style={circle,minimum size=7mm,inner sep=0pt,text=white,font=\bfseries\scriptsize}]
  \node[ch,fill=attncol] at (0,0) {$+$}; \node[below,font=\scriptsize] at (0,-0.45) {$A$};
  \node[ch,fill=attncol] at (2.4,0) {$+$}; \node[below,font=\scriptsize] at (2.4,-0.45) {$B$};
  \node[ch,fill=defcol] at (4.8,0) {$-$}; \node[below,font=\scriptsize] at (4.8,-0.45) {$C$};
  \draw[<->,black!45,dashed] (0,0.55) -- (2.4,0.55) node[midway,above,font=\scriptsize]{$\SI{10}{\centi\metre}$};
  \draw[<->,black!45,dashed] (2.4,0.55) -- (4.8,0.55) node[midway,above,font=\scriptsize]{$\SI{10}{\centi\metre}$};
\end{tikzpicture}
\end{center}
\begin{enumerate}[label=\alph*)]
  \item Donne l'intensité et le sens de la force exercée sur $B$ par $A$, puis par $C$.
  \item En déduire l'intensité et le sens de la force totale subie par $B$.
\end{enumerate}
\end{exo}

\begin{exo}[superposition à angle droit : calcul complet]
Une charge $q_0 = \SI{+2}{\micro\coulomb}$ est placée en $O$. Une charge
$q_1 = \SI{+2}{\micro\coulomb}$ est à $\SI{20}{\centi\metre}$ sur l'axe horizontal, une charge
$q_2 = \SI{+6}{\micro\coulomb}$ est à $\SI{30}{\centi\metre}$ sur l'axe vertical.
\begin{center}
\begin{tikzpicture}[scale=0.95,>=Stealth,font=\footnotesize,
    ch/.style={circle,fill=attncol,minimum size=7mm,inner sep=0pt,text=white,font=\bfseries\scriptsize}]
  \node[ch] (O) at (0,0) {$q_0$};
  \node[ch] (q1) at (3,0) {$q_1$};
  \node[ch,minimum size=8.5mm] (q2) at (0,2.6) {$q_2$};
  \draw[<->,black!45,dashed] (0,-0.55) -- (3,-0.55) node[midway,below,font=\scriptsize]{$\SI{20}{\centi\metre}$};
  \draw[<->,black!45,dashed] (-0.6,0) -- (-0.6,2.6) node[midway,left,font=\scriptsize]{$\SI{30}{\centi\metre}$};
  \draw[->,line width=1.1pt,defcol] (O) -- (-1.3,0) node[left,font=\scriptsize]{$\vv{F}_1$};
  \draw[->,line width=1.1pt,defcol] (O) -- (0,-1.3) node[below,font=\scriptsize]{$\vv{F}_2$};
\end{tikzpicture}
\end{center}
\begin{enumerate}[label=\alph*)]
  \item Calcule les intensités $F_1$ (due à $q_1$) et $F_2$ (due à $q_2$) subies par $q_0$.
  \item En déduire l'intensité de la force résultante sur $q_0$.
\end{enumerate}
\end{exo}

\end{document}
